\subsection{Measurable mappings and scalarly measurable mappings}

If a mapping f of T into a Hausdorff locally convex space F is scalarly \(\mu\) -measurable, it does not in general follow that f is \(\mu\) -measurable (Exer. 12). Nevertheless:

PROPOSITION 12. --- If F is a separable, metrizable locally convex space, then every scalarly \(\mu\) -measurable mapping f of T into F is \(\mu\) -measurable.

For, F may be regarded as a subspace of a countable product \(\prod_{n}E_{n}\) of Banach spaces (TVS, II, §1, No.~3, Prop. 3), and one can suppose that \(\mathrm{pr}_n(\mathbf{F})\) is dense in \(\mathbf{E}_n\) , which is therefore separable. For every \(n\) , the mapping \(\mathrm{pr}_n \circ \mathbf{f}\) is scalarly \(\mu\) -measurable, hence \(\mu\) -measurable (Ch. IV, §5, No.~5, Cor. 2 of Prop. 10), consequently \(\mathbf{f}\) is \(\mu\) -measurable (Ch. IV, §5, No.~3, Th. 1).

PROPOSITION 13. --- Let F be a locally convex space that is the direct limit of a sequence of separable, metrizable locally convex spaces \(F_{n}\) , F being the union of the \(F_{n}\) . Let \(F'\) be the dual of F, equipped with the topology \(\sigma(\mathbf{F}',\mathbf{F})\) . Then, every scalarly \(\mu\) -measurable mapping f of T into \(F'\) is \(\mu\) -measurable.

Suppose first that F is metrizable and separable, and let D be a countable dense set in F. Let \((\mathrm{V}_{n})\) be a decreasing fundamental sequence of balanced, convex, open neighborhoods of 0 in F; the polar sets \(V_{n}^{\circ}\) are equicontinuous and their union is all of \(F'\) . Let \(\mathrm{T}_{n} = \mathbf{f}^{-1}(\mathrm{V}_{n}^{\circ})\) ; the sequence \((\mathrm{T}_{n})\) is increasing and \(T = \bigcup_{n} T_{n}\) ; let us show that each of the \(T_{n}\) is \(\mu\) -measurable. Indeed, \(D \cap V_{n}\) is dense in \(V_{n}\) ; for every \(y \in D \cap V_{n}\) , let \(S_{y}\) be the set of \(t \in T\) such that \(|\langle y, f(t) \rangle| \leqslant 1\) ; the hypothesis implies that each of the \(S_{y}\) is measurable, and \(T_{n}\) is the intersection of the countable family of the \(S_{y}\) \((y \in D \cap V_{n})\) . This being so, for every compact subset K of T and every \(\varepsilon > 0\) , there exists an integer n such that \(\mu\left(\mathrm{K}-\left(\mathrm{K}\cap\mathrm{T}_{n}\right)\right)\leqslant\frac{\varepsilon}{4}\) , then a compact subset \(K_{1}\) of \(K\cap T_{n}\) such that \(\mu\left(\left(\mathrm{K}\cap\mathrm{T}_{n}\right)-\mathrm{K}_{1}\right)\leqslant\frac{\varepsilon}{4}\) ; finally, there exists a compact subset \(K_{2}\) of \(K_{1}\) such that \(\mu(K_{1}-K_{2})\leqslant\frac{\varepsilon}{2}\) and such that the restrictions to \(K_{2}\) of all of the functions \(\langle y,f\rangle\) , where \(y\in D\) , are continuous (Ch. IV, §5, No.~1, Prop. 2). Since the set \(\mathbf{f}(K_{2})\subset\mathbf{f}(T_{n})\subset V_{n}^{\circ}\) is equicontinuous, the topology induced by \(\sigma(F',F)\) on \(\mathbf{f}(K_{2})\) is identical to the topology of pointwise convergence in D (GT, X, §2, No.~4, Th. 1); consequently, the restriction of f to \(K_{2}\) is continuous, whence our assertion in the first case.

Let us pass to the general case. If \(z^{\prime}\) is a continuous linear form on F, its restriction \(z_{n}^{\prime}\) to \(F_{n}\) is continuous; since \(F = \bigcup_{n} F_{n}\) , the dual \(F^{\prime}\) of F may be identified (algebraically) with a linear subspace of the product \(\prod_{n} F_{n}^{\prime}\) , and \(pr_{n}z^{\prime} = z_{n}^{\prime}\) . Moreover, since each finite subset of F is contained in one of the \(F_{n}\) , the topology \(\sigma(F^{\prime}, F)\) is none other than the topology induced by the product topology of the topologies \(\sigma(F_{n}^{\prime}, F_{n})\) . This being so, if f is scalarly \(\mu\) -measurable then \(pr_{n} \circ f\) is scalarly \(\mu\) -measurable, since for every \(t \in T\) , \(\operatorname{pr}_{n}(f(t))\) is the restriction of \(f(t)\) to \(F_{n}\) . The first part of the proof shows that \(pr_{n} \circ f\) is \(\mu\) -measurable for every n, therefore so is f (Ch. IV, §5, No.~3, Th. 1).

